\section{Introduction}
Test generation requires a connection between what software should do and observations that can expose when it does something else. Constructing that connection involves several kinds of work: interpreting requirements, selecting inputs, arranging an execution environment, choosing observations, and expressing an expected result. Large language models (LLMs) can contribute to each of these activities because they operate across natural language, program code, and semi-structured technical documentation. Research has consequently expanded from generating unit-test code to producing acceptance scenarios, postconditions, API call sequences, and fuzzing drivers~\cite{S01,S02,S05}.

This breadth creates a comparison problem. Two approaches can both generate executable tests while having different access to the system under test (\sut), different sources of expected behavior, and different intended uses. A regression generator may inspect an implementation and retain tests that pass on it. A requirements-driven generator may need to produce an oracle before an implementation exists. An issue-reproduction agent may receive a report of an already observed failure. The same passing test carries a different meaning in each setting. Comparing their coverage or execution success without describing these conditions obscures what has actually been learned about test quality.

Recent systems make this distinction especially important. TestGen-LLM improves existing tests through operational filters, including execution and coverage checks~\cite{U07}. CoverUp uses coverage feedback to guide Python test generation~\cite{U05}. In a different setting, PyTester learns to produce tests from textual descriptions, using executable reference information during training~\cite{T01}. UTRL couples test and code generation through adversarial reinforcement learning~\cite{T02}. These systems differ in where behavior is observed, when it is available, and how the observations influence the final test. Their common use of an LLM does not determine their testing objective.

Existing surveys provide substantial foundations. Broad reviews map LLM applications across testing and repair~\cite{S01,S04}; unit-testing surveys organize techniques, model adaptation, and quality assurance~\cite{S02,S03}; requirements-based reviews cover transformations from natural-language and formal specifications~\cite{S05,S06}. A recent systematic review directly addresses LLM-based test-case generation~\cite{S07}. Oracle generation, hybrid techniques, and evaluation limitations already appear in these accounts. Our organizing question is how these aspects interact: \emph{what information produces a test, and what evidence justifies using that test to judge software?}

We develop this question through a structured narrative synthesis of a curated collection of 60 source records. The collection includes seven surveys and roadmaps, together with primary methods, benchmarks, empirical studies, and selected neural predecessors. Its breadth supports a cross-setting analysis; its preliminary screening status is described in Section~\ref{sec:scope}. We make three contributions. First, we define a framework that separates testing objectives, generated artifacts, information access, and quality evidence. Second, we use the framework to connect prompting, program analysis, feedback, agents, and model training across testing settings. Third, we organize evaluation designs according to the claims they can support and derive concrete research questions from the remaining mismatches.

The central argument is that test generation should be understood as both an artifact-construction problem and an evidence-construction problem. Better syntax and broader execution improve the first. Confidence in expected behavior requires the second. This distinction helps explain when tools from different research traditions can be combined and which claims require additional validation.
